\chapter{绪论}

\section{研究背景与意义}

\subsection{天地一体化网络与低轨卫星发展}

% TODO: 天地一体化网络是新一代移动网络的发展方向，LEO卫星（如Starlink）的爆发式增长

\subsection{非法卫星终端的安全威胁}

% TODO: 非法终端利用卫星资源的安全隐患，定位需求的紧迫性

\subsection{低成本无人机定位平台的优势}

% TODO: 无人机的灵活性、低成本、易部署特点

\section{国内外研究现状}

\subsection{无线电定位技术研究现状}

% TODO: TOA/TDOA/AOA/RSSI等传统几何定位技术，各自优缺点

\subsection{射线追踪与指纹定位研究现状}

% TODO: 射线追踪信道建模、RF指纹定位（CNN/ResNet等深度学习方法）

\subsection{多模态信息融合定位研究现状}

% TODO: EKF/粒子滤波等传统方法、基于强化学习的融合方法、物理感知缺失的批判

\subsection{多目标定位与检测技术研究现状}

% TODO: MUSIC/ESPRIT谱估计、基于深度学习的目标检测、YOLO系列发展

\section{现有研究的不足}

% TODO: 三大痛点总结：
% 1. 单一模态在城市NLOS环境失效
% 2. 传统融合方法缺乏环境自适应性（重点批判"物理感知缺失"）
% 3. 多目标串行算法的计算瓶颈与虚警问题

\section{主要研究内容与创新点}

\subsection{研究内容}

% TODO: 概述两大研究方向
% 研究内容一：基于射线追踪的信道建模 + ResNet指纹定位 + PA-DQN融合
% 研究内容二：信号成像Radio Map + Anchor-Free YOLOv7多目标检测

\subsection{主要创新点}

% TODO: 三大创新点
% 创新点一：射线追踪 + ResNet-50深度RF指纹 + MIMO空间分集 → 复杂环境精准定位
% 创新点二：PA-DQN物理感知RL融合框架 → 动态自适应多模态融合
% 创新点三：信号成像 + Anchor-Free YOLOv7 → 跨模态多目标并行定位

\section{论文组织结构}

% TODO: 各章内容简述

\section{本章小结}

% TODO: 简要小结本章内容
